\section{Không có tầng đáy}
\label{sec:ch11-khong-co-tang-day}

\index{độ trung thực!chỗ đứng trong bài báo bậc hai|(}%
Không cái nền nào trong ba cái ấy là nền của độ trung thực, thế mà chữ ấy vẫn
có mặt trong bài. Đọc xem nó có mặt ở đâu là cách gọn nhất để thấy chỗ đứng
thật của nó.

Trong cả 32 trang, chữ \enquote{faithful} và các dạng của nó xuất hiện mười ba
lần. Sáu lần trong số đó nằm bên trong tên riêng của một phương pháp khác,
faithful Shapley interaction index của Tsai và cộng sự, nên chúng không nói gì
về tính chất; giữ nguyên tên tiếng Anh ở đây là cần thiết, vì chính chuỗi chữ
ấy là thứ đang được đếm. Hai lần là tiêu chuẩn phân tích của
mục~\ref{sec:ch11-ro-ri-tuong-tac}. Một lần nói rằng AttnLRP đạt mức trung thực
tốt nhất hiện có trong các bài toán ngôn ngữ và thị giác, và câu ấy quy về một
bài báo khác chứ không về phép đo nào của bài này. Một lần nữa nói về cuộc
tranh luận cơ chế chú ý có giải thích trung thực transformer hay không, tức về
tài liệu khác. Hai lần còn lại nằm ở phần thí nghiệm thứ hai và nói về chính
bài: một lần gọi biến thể Meta-MaskSigLIP là trung thực nhất để diễn giải
SigLIP-2, một lần nói tới việc gán trung thực các tương tác giữa hai token. Cả
hai đều dựa trên chỉ số nhận diện tương tác mà mục trước đã đọc, tức một phép
đo trên nhãn của bộ dữ liệu~\autocite{p23metagame}.

Lần thứ mười ba là lần quan trọng, và nó nằm trong phần hạn chế của bài. Các
tác giả viết rằng phần phân tích trên mô hình ngôn ngữ chỉ là một minh họa trên
vài lời nhắc, và rằng một phép so sánh phương pháp cho ra kết luận chặt thì còn
cần thêm các phép đo độ trung thực, chẳng hạn các đường cong deletion và
insertion trên từng cặp token~\autocite{p23metagame}.

Câu ấy là chỗ chương này gặp lại chương~\ref{ch:do-do-trung-thuc}. Deletion và
insertion là hai trong số các họ chỉ số mà mục~\ref{sec:ch08-ho-xoa-dan} đã đọc
kỹ, và điều chương ấy kết luận về cả bốn họ là chưa có gì xác lập rằng chúng đo
đúng đại lượng ghi trên tên chúng. Cho nên bước tiếp theo mà bài tự nêu ra cho
mình dẫn thẳng tới một dụng cụ chưa được kiểm định. Không phải một dụng cụ đã
biết là hỏng, vì mục~\ref{sec:ch09-doi-tuong-khac-nhau} đã dựng ranh giới chặn
đúng phép suy rộng ấy; là một dụng cụ mà chưa ai đem đối chiếu với ground truth
nào.

\begin{figure}[tbp]
  \centering
  \begin{tikzpicture}[
  every node/.style={font=\footnotesize, align=center},
  obj/.style={draw=black!65, rounded corners=2pt, fill=black!5,
    minimum width=34mm, minimum height=11mm},
  ok/.style={draw=black!75, rounded corners=2pt, fill=black!10,
    minimum width=40mm, minimum height=13mm},
  gap/.style={draw=black!45, dotted, thick, rounded corners=2pt, fill=white,
    minimum width=40mm, minimum height=13mm},
  up/.style={-{Stealth[length=2mm]}, thick, black!70},
  chk/.style={-{Stealth[length=2mm]}, thick, black!70},
  miss/.style={-{Stealth[length=2mm]}, thick, black!45, dotted},
  lbl/.style={font=\scriptsize, align=center, black!55}
]
  \node[obj] (model)  at (0,3.0) {Mô hình $f$};
  \node[obj] (attr)   at (0,0.6) {Attribution bậc nhất\\$\varphi_i(x,f)$};
  \node[obj] (meta)   at (0,-1.8) {Meta-attribution\\$\varphi_{j\to i}(x,f)$};

  \node[ok]  (acc)    at (6.4,3.0) {Độ chính xác\\trên tập kiểm định};
  \node[gap] (metric) at (6.4,0.6) {Chỉ số trung thực:\\chưa được kiểm định};
  \node[gap] (proxy)  at (6.4,-1.8) {Nhãn của người và của\\bộ dữ liệu: tính hợp lý};

  \draw[up] (attr) -- node[lbl, right=1mm] {giải thích} (model);
  \draw[up] (meta) -- node[lbl, right=1mm] {giải thích} (attr);

  \draw[chk]  (acc)    -- (model);
  \draw[miss] (metric) -- (attr);
  \draw[miss] (proxy)  -- (meta);

  \draw[miss] (proxy.east) -- ++(1.0,0) |- (metric.east);
  \node[lbl, text width=26mm] at (9.1,-0.6)
    {bước bài báo nêu\\mà chưa làm};
\end{tikzpicture}

  \caption{Ba tầng của chương này và cái được dùng để kiểm chứng từng tầng. Nét
  chấm đánh dấu tầng không có dụng cụ đo độ trung thực đã được kiểm định, vì hai
  lý do khác nhau: tầng hai có dụng cụ nhưng dụng cụ ấy chưa được kiểm định,
  tầng ba có phép đo nhưng phép đo ấy đo thứ khác. Mũi tên chấm bên phải là
  bước mà phần hạn chế của bài nêu ra và chưa thực hiện.}
  \label{fig:ch11-ba-tang}
\end{figure}

Hình~\ref{fig:ch11-ba-tang} xếp ba tầng ấy lại. Tầng trên cùng là mô hình, và
tầng ấy có một dụng cụ đo mà cả lĩnh vực chấp nhận: độ chính xác trên tập kiểm
định, tính được, so được giữa hai mô hình, và không cần lời giải thích nào. Tầng
thứ hai là attribution bậc nhất, và dụng cụ của nó là các họ \tn{chỉ số trung thực}{faithfulness metric}
của chương~\ref{ch:do-do-trung-thuc}, chưa cái nào được kiểm định. Tầng thứ ba
là meta-attribution, và cái đứng ở chỗ dụng cụ của nó là ba phép đo của
mục~\ref{sec:ch11-ba-phep-do-va-ba-cai-nen}, cả ba đều nằm ở phía tính hợp lý.
Mũi tên chấm nối tầng ba ngược lên tầng hai là câu trong phần hạn chế: bước
đúng đắn tiếp theo mà bài nêu ra là mượn dụng cụ của tầng hai, tức mượn đúng
cái đang thiếu kiểm định.

Đó là nghĩa thứ hai của nhan đề, và là nghĩa mà sách đọc nó theo. Mỗi tầng được
giải thích bằng một tầng nữa bên dưới, còn câu hỏi tầng này có đúng hay không
thì được đẩy xuống tầng dưới, chỗ nó đã không có câu trả lời từ trước.

Không nên đọc bài báo này nặng hơn mức nó đáng chịu. Các tác giả không giấu
khoảng trống ấy; chính họ viết nó ra, gọi tên phép đo còn thiếu, và chỉ ra dụng
cụ họ sẽ dùng. Đó là lý do bài này dùng được làm bằng chứng cho
lập luận của chương: chỗ hở được người trong cuộc ghi lại chứ không phải do
sách phát hiện. Điều sách thêm vào là một quan sát mà bài không có chỗ để nêu,
rằng dụng cụ được chỉ ra ấy nằm đúng trong tập dụng cụ mà Phần~IV đang chất
vấn.

Chữ trong nhan đề còn một chỗ dễ nhầm nữa. Trong tài liệu về đánh giá XAI có
một khái niệm gần âm là bài toán siêu đánh giá, mà bài mô tả bằng đúng một vế
câu: nó chấm chất lượng của attribution. Bài tự tách mình khỏi khái niệm ấy
trong phần công trình liên quan, ghi rằng metagame chỉ trùng tên, còn đối tượng
của metagame là hiệu ứng \tn{tương tác bậc hai}{second-order
interaction}~\autocite{p23metagame}.
Sách chưa đọc bài về siêu đánh giá, nên không thêm khóa trích dẫn cho nó ở đây
và không nói gì thêm về nội dung của nó;
chương~\ref{ch:khoang-trong-nghien-cuu} sẽ phải quyết định có đọc hay không.
\index{độ trung thực!chỗ đứng trong bài báo bậc hai|)}%
